% 矮行星（综述）
% license CCBYSA3
% type Wiki

本文根据 CC-BY-SA 协议转载翻译自维基百科\href{https://en.wikipedia.org/wiki/Dwarf_planet}{相关文章}。

\begin{figure}[ht]
\centering
\includegraphics[width=6cm]{./figures/d21611a5210b90a2.png}
\caption{} \label{fig_Dwarf_1}
\end{figure}
“矮行星是一类直接绕太阳运行、具有行星质量但体积较小的天体。它们的质量足以使自身达到由重力塑形的近球形，但不足以像太阳系中八大经典行星那样在轨道上实现支配地位。矮行星的原型是冥王星，它曾被视为一颗行星长达数十年，直到 2006 年‘矮行星’概念被正式采用。许多行星地质学家认为矮行星与具有行星质量的卫星都应被视为行星[1]，但自 2006 年以来，国际天文学联合会（IAU）和许多天文学家已将它们排除在行星名录之外。

矮行星可能具有地质活动能力——这一预期在 2015 年“黎明号”对谷神星的探测和“新视野号”对冥王星的探测中得到证实。因此，它们成为行星地质学家特别关注的对象。

天文学家普遍同意至少有九个最大的候选体属于矮行星——按直径从大到小大致排列为：冥王星、厄里斯、妊神星、鸟神星、公公星、夸欧尔、塞德娜、谷神星和奥库斯。对于第十大候选体萨拉西娅仍存在相当的不确定性，因此可将其视为临界案例。在这十个天体中，已有两个被航天器访问过（冥王星和谷神星），另有七个至少拥有一颗已知卫星（厄里斯、妊神星、鸟神星、公公星、夸欧尔、奥库斯和萨拉西娅），这使得它们的质量得以确定，从而可估算其密度。质量和密度又可用于构建地球物理模型，以尝试判断这些世界的性质。只有塞德娜既未被探测器访问，也没有已知卫星，因此其质量难以准确估计。一些天文学家还将许多更小的天体也纳入矮行星的范畴[2]，但目前并无共识认为这些小天体很可能是矮行星。”**
\subsection{概念的历史}
\begin{figure}[ht]
\centering
\includegraphics[width=6cm]{./figures/0ce0512cdfe2e0c4.png}
\caption{} \label{fig_Dwarf_2}
\end{figure}
\begin{figure}[ht]
\centering
\includegraphics[width=6cm]{./figures/6ae4d90716be2481.png}
\caption{} \label{fig_Dwarf_3}
\end{figure}
自1801年起，天文学家发现了谷神星以及位于火星与木星之间的其他天体，这些天体在随后的数十年间都被视为行星。自那时直到大约1851年，当行星数量已达到23颗时，天文学家开始使用“asteroid”（来自希腊语，意为“类星的”或“星形的”）一词来指代较小的天体，并开始将它们区分为小行星而非大行星。[4]

随着1930年冥王星的发现，大多数天文学家认为太阳系包含九颗主要行星，以及成千上万的显著更小的天体（小行星与彗星）。在将近50年的时间里，人们曾认为冥王星比水星更大，[5][6] 但在1978年冥王星的卫星卡戎被发现后，天文学家得以精确测量冥王星的质量，并确定其远小于最初的估计。[7] 其质量大约只有水星的二十分之一，使冥王星成为迄今最小的行星。尽管冥王星仍比小行星带中最大的天体谷神星大十倍以上，但其质量却只有地球月球的五分之一。[8] 此外，由于其具有一些异常特征，如较大的轨道离心率与较高的轨道倾角，愈发明显冥王星与其他行星类型不同。[9]

在1990年代，天文学家开始在与冥王星相同的空间区域（即如今所称的柯伊伯带）以及更远处发现更多天体。[10] 其中许多与冥王星共享关键的轨道特征，冥王星开始被视为一类新天体的最大成员，即冥族（plutinos）。于是变得明显：要么这些较大的天体也需要被归类为行星，要么冥王星需要像谷神星在更多小行星被发现后那样被重新分类。[11] 这使得一些天文学家不再将冥王星称为行星。多个术语包括“subplanet”（次行星）与“planetoid”（类行星体）等，被用于指代如今所谓的矮行星。[12][13] 天文学家也确信会有更多与冥王星大小相当的天体被发现，如果冥王星继续被视为行星，行星数量将会快速增长。[14]

海王星外天体厄里斯（当时称为 2003 UB313）于 2005 年 1 月被发现；[15] 当时认为其尺寸略大于冥王星，因此一些报道非正式地称其为“第十颗行星”。[16] 因此，该问题在 2006 年 8 月的国际天文学联合会（IAU）大会上成为激烈争论的焦点。[17] IAU 的最初草案提议将卡戎、厄里斯与谷神星列入行星名单。由于许多天文学家反对这一方案，乌拉圭天文学家 Julio Ángel Fernández 与 Gonzalo Tancredi 制定了另一项替代方案：他们提出设立一个中间类别，用于归类那些质量足以达到流体静力平衡（呈球形），但并未清除其轨道中微行星的天体。除了将卡戎从名单中移除之外，该新方案也将冥王星、谷神星与厄里斯排除在行星名单之外，因为它们均未清除其所在轨道。[18]

尽管有人对此分类在系外行星中的适用性提出了担忧，[19] 该问题仍未得到解决；提议是等到能够观测到与矮行星体积相当的系外天体时再作决定。[18]

在 IAU 给出“矮行星”定义的直接后续中，一些科学家表达了他们对 IAU 决议的不满。[20] 抗议活动包括汽车保险杠贴纸与 T 恤标语。[21] 厄里斯的发现者 Mike Brown 则赞同将行星数量减少至八颗的做法。[22]

NASA 在 2006 年宣布将采用 IAU 制定的新指南。[23] NASA 冥王星任务的负责人 Alan Stern 不认可当前 IAU 的行星定义，无论是将矮行星排除在行星类别之外，还是依据轨道特征（而非天体的固有特征）来定义矮行星。[24] 因此，他在 2011 年仍将冥王星称为行星，[25] 并将其他可能的矮行星，如谷神星与厄里斯，以及更大的卫星，也视作额外的行星。[26] 在 IAU 发布定义的数年前，他便使用轨道特征将“überplanets”（主导性的八颗行星）与“unterplanets”（矮行星）区分开来，并认为二者皆为“行星”。[27]
\subsection{名称}

